\chapter{Revisão da Literatura}
\label{chap:revisao_literatura}

É fato que a NFV traz, por sua natureza, flexibilidade sobre o uso das VNFs. Por outro lado, surge a necessidade de processar decisões sobre as atividades de implantação e até migração de serviços virtualizados. Para que esse processo ocorra, é necessário compreender como as atividades de monitoramento se integram aos componentes previstos para uma arquitetura NFV. Assim, este Capítulo apresenta os fundamentos da NFV e o ciclo de vida das VNFs (Seção~\ref{sec:nfv}), seguidos das abordagens de monitoramento no espaço do usuário Sysstat e Prometheus (Seção~\ref{sec:sysstat_prometheus}), da tecnologia eBPF e seu modelo de programabilidade no \textit{kernel} (Seção~\ref{sec:ebpf}), e de uma análise comparativa entre instrumentação via \textit{kprobes} e leitura de \texttt{/proc} (Seção~\ref{sec:kprobes_proc}), concluindo com os trabalhos relacionados que contextualizam a contribuição deste estudo (Seção~\ref{sec:trabalhos_relacionados}).

\section{Virtualização de Funções de Rede}
\label{sec:nfv}

A NFV é um paradigma arquitetural formalizado pelo Instituto Europeu de Normas de Telecomunicações (\textit{European Telecommunications Standards Institute} - ETSI), que propõe desacoplar as funções de rede de equipamentos proprietários, permitindo que sejam executadas como software em servidores de propósito geral \cite{etsi_nfv_002_2014, etsi_nfv_006_2024, mijumbi2016}. Essa abordagem elimina a dependência de hardware especializado, reduz custos operacionais e aumenta a flexibilidade no provisionamento de serviços \cite{7045396, 10.1145/3397022}. VNFs comuns incluem roteadores virtuais, Sistemas de Detecção de Intrusão (\textit{Intrusion Detection Systems} - IDS), balanceadores de carga e \textit{Firewalls} de Aplicação Web (\textit{Web Application Firewall} - WAF), entre outros \cite{10.1145/3172866}.

A arquitetura de referência definida pelo ETSI organiza os componentes NFV em três domínios funcionais principais \cite{etsi_nfv_002_2014, etsi_nfv_006_2024, mijumbi2016, repec:wly:intnem:v:31:y:2021:i:5:n:e2068} descritos a seguir:

\begin{itemize}
    \item \textbf{NFVI (\textit{NFV Infrastructure})}: camada que engloba os recursos físicos de computação, armazenamento e rede, além da camada de virtualização que os abstrai em recursos virtuais disponíveis para as VNFs;
    \item \textbf{VNFs (\textit{Virtualized Network Functions})}: funções de rede implementadas como software, executando sobre os recursos da NFVI, interconectadas por redes virtuais e gerenciadas por descritores padronizados;
    \item \textbf{MANO (\textit{Management and Orchestration})}: bloco responsável pelo gerenciamento e orquestração de toda a infraestrutura NFV, composto pelo \textit{NFV Orchestrator} (NFVO), que gerencia os serviços de rede; pelo \textit{VNF Manager} (VNFM), que controla o ciclo de vida das VNFs; e pelo \textit{Virtualized Infrastructure Manager} (VIM), que administra os recursos computacionais da NFVI.
\end{itemize}

\subsection{Ciclo de Vida das VNFs}

O gerenciamento do ciclo de vida de uma VNF compreende as operações de instanciação, escalonamento (\textit{scale out/in}), recuperação (\textit{healing}), atualização e término \cite{repec:wly:intnem:v:31:y:2021:i:5:n:e2068}. Essas operações são coordenadas pelo VNFM e dependem de informações precisas sobre o estado operacional da VNF, incluindo consumo de CPU, memória e tráfego de rede, para decisões de escalonamento automático e recuperação de falhas \cite{10.23919}. A qualidade e a latência dessas informações determinam diretamente a eficiência das operações de gerenciamento: métricas coletadas com atraso elevado comprometem a capacidade de reação do MANO a variações abruptas de carga \cite{oliveira_observability_2026}.

Para ilustrar a dinâmica dessas operações, considere alguns cenários típicos. Em um cenário de escalonamento horizontal (\textit{scale out}), um IDS que monitora tráfego de entrada pode atingir o limite de capacidade de processamento durante um pico de acessos; o VNFM, ao detectar o consumo elevado de CPU por meio das métricas coletadas, instancia automaticamente uma nova réplica da VNF e redistribui o tráfego entre as instâncias \cite{NIKBAZM2022102632, repec:wly:intnem:v:31:y:2021:i:5:n:e2068}. Em um cenário de recuperação (\textit{healing}), uma instância de WAF que deixa de responder a verificações de integridade (\textit{health checks}) é automaticamente encerrada e substituída por uma nova instância provisionada a partir do mesmo descritor, restaurando o serviço sem intervenção manual \cite{repec:wly:intnem:v:31:y:2021:i:5:n:e2068, mijumbi2016}.

Já em um cenário de atualização (\textit{update}), uma nova versão de um balanceador de carga virtual pode ser implantada de forma gradual (\textit{rolling update}), substituindo instâncias antigas enquanto o tráfego é progressivamente redirecionado, minimizando a interrupção do serviço \cite{repec:wly:intnem:v:31:y:2021:i:5:n:e2068}. Em todos esses cenários, a precisão e a latência das métricas de monitoramento são determinantes: uma leitura atrasada ou imprecisa do estado da VNF pode levar o MANO a tomar decisões equivocadas, como instanciar réplicas desnecessárias ou não detectar falhas a tempo \cite{10.1145/3397022, oliveira_observability_2026}.

O monitoramento individual de cada VNF da cadeia é indispensável para isolar gargalos, atribuir responsabilidade pelo atraso introduzido em cada elo e detectar anomalias de forma antecipada \cite{10.23919, oliveira_observability_2026}. No contexto deste trabalho, o WAF representa um elo isolado dessa cadeia, funcionando como VNF representativa para quantificar a sobrecarga de diferentes abordagens de monitoramento. A extensão das abordagens avaliadas para o monitoramento de SFCs completas, considerando a propagação de sobrecarga ao longo de múltiplos elos encadeados, representa uma direção promissora para trabalhos futuros.

O ciclo de vida de VNFs pode se tornar mais complexo quando o Encadeamento de Funções de Serviço (\textit{Service Function Chaining} - SFC) se apresenta como requisito para o ambiente. SFCs definem sequências ordenadas de VNFs em que o tráfego de rede deve percorrer para que um serviço seja provisionado corretamente \cite{rfc7665, 10.1145/3172866}. Em uma cadeia típica, os pacotes podem passar por um \textit{firewall}, em seguida por um IDS e finalmente por um balanceador de carga antes de atingirem o servidor de destino. O monitoramento individual de cada VNF da cadeia é indispensável para isolar gargalos, atribuir responsabilidade pelo atraso introduzido em cada elo e detectar anomalias de forma antecipada \cite{10.23919, oliveira_observability_2026}. No contexto deste trabalho, o WAF representa um elo isolado dessa cadeia, funcionando como VNF representativa para quantificar a sobrecarga de diferentes abordagens de monitoramento.

\subsection{Requisitos de Monitoramento em Ambientes NFV}
\label{subsec:requisitos_monitoramento}

O monitoramento de VNFs não se limita à coleta de métricas de infraestrutura: envolve também a observação do comportamento funcional das funções virtualizadas sob diferentes cargas de tráfego \cite{10.23919}. Nesse contexto, \citeonline{oliveira_observability_2026} argumentam que a observabilidade constitui o fator ausente no gerenciamento de ambientes NFV, pois as ferramentas de orquestração dependem de telemetria precisa e de baixa latência para manter os Acordos de Nível de Serviço (\textit{Service Level Agreement} - SLA). A sobrecarga gerada pelas camadas de virtualização e pelo tráfego intenso exige mecanismos de monitoramento altamente eficientes, uma vez que ferramentas intrusivas competem diretamente pelos recursos computacionais destinados ao processamento do tráfego \cite{10.1145/3397022, 10.1109/NOMS47738.2020.9110434}.

Para ilustrar a importância do monitoramento, considere um cenário de segurança em que um WAF virtualizado protege uma aplicação Web de alto tráfego. Durante um ataque Distribuído de Negação de Serviço (\textit{Distributed Denial of Service} - DDoS), o volume de requisições pode aumentar abruptamente em uma ou duas ordens de grandeza em questão de segundos \cite{technologies12080122}. Nesse contexto, o MANO depende inteiramente das métricas coletadas pelo monitor para detectar a saturação da VNF e acionar o escalonamento horizontal a tempo de evitar violação de SLA. Se o mecanismo de monitoramento introduz latência de coleta superior ao intervalo de reação esperado pelo orquestrador, a decisão de escalonamento chega tarde, e requisições legítimas são descartadas ou atrasadas além do limiar aceitável \cite{oliveira_observability_2026, 10.23919}.
Um segundo cenário ilustrativo envolve o monitoramento de um IDS virtualizado em um ambiente de borda (\textit{edge computing}). Nesse caso, os recursos computacionais disponíveis são tipicamente escassos, e qualquer sobrecarga introduzida pela ferramenta de monitoramento reduz diretamente a capacidade de inspeção de pacotes da VNF \cite{10.1109/NOMS47738.2020.9110434}.

Estudos demonstram que ferramentas baseadas em \textit{polling} periódico via \texttt{/proc} podem consumir entre 3\% e 8\% de um núcleo de CPU sob alta frequência de amostragem, o que em ambientes de borda representa uma fração não negligenciável da capacidade total disponível \cite{10.1145/3452296.3472888}. Esse efeito é agravado quando múltiplas VNFs são monitoradas simultaneamente no mesmo ambiente, pois a sobrecarga de cada agente se acumula de forma aditiva.
Um terceiro cenário relevante diz respeito à detecção de anomalias funcionais em tempo real. Em um balanceador de carga virtualizado, picos abruptos de latência em conexões TCP individuais podem indicar falhas silenciosas em servidores de \textit{backend}, que não se manifestam imediatamente como indisponibilidade total da VNF \cite{epping2024rtt}.

Ferramentas de monitoramento baseadas em leituras periódicas de contadores agregados, como as disponíveis via \texttt{/proc/net/dev}, são incapazes de capturar esses eventos transitórios, pois calculam médias sobre o intervalo de amostragem e mascaram variações de curta duração \cite{godard2024sysstat}. Em contraste, abordagens baseadas em instrumentação direta no \textit{kernel}, como \textit{kprobes} associadas a funções TCP, permitem registrar a latência de cada conexão individualmente, viabilizando a detecção de anomalias com granularidade de microssegundos \cite{10.1109/NOMS47738.2020.9110434, epping2024rtt}.
Além da granularidade temporal, os requisitos de monitoramento em ambientes NFV incluem a capacidade de atribuir métricas a fluxos ou conexões específicas, e não apenas a interfaces de rede agregadas \cite{rfc9232}.

Em um SFC composto por \textit{firewalls}, IDS e balanceadores de carga, por exemplo, é necessário identificar em qual elo da cadeia um determinado fluxo sofreu atraso ou foi descartado, o que exige correlação entre métricas coletadas em diferentes VNFs \cite{10.23919, oliveira_observability_2026}. Essa capacidade de correlação depende diretamente da precisão e da resolução temporal das métricas coletadas em cada elo, reforçando a necessidade de mecanismos de monitoramento de baixa intrusividade e alta granularidade.
Por fim, em ambientes de produção com alta densidade de VNFs, o próprio tráfego de telemetria gerado pelos agentes de monitoramento pode competir com o tráfego de serviço pelos recursos de rede e CPU do host \cite{kannan2024designing}. Esse fenômeno, conhecido como \textit{monitoring overhead amplification}, é particularmente crítico quando múltiplos exporters transmitem métricas simultaneamente para um servidor central de coleta, como no modelo \textit{pull} do Prometheus \cite{usman_observability_2022}. A escolha da abordagem de monitoramento, portanto, não é apenas uma decisão técnica de implementação, mas um fator determinante para a qualidade de serviço oferecida pelas VNFs monitoradas.

\section{Monitoramento no Espaço do Usuário: Sysstat e Prometheus}
\label{sec:sysstat_prometheus}

As ferramentas tradicionais de monitoramento operam majoritariamente no espaço do usuário (\textit{user-space}), extraindo dados por meio do pseudossistema de arquivos \texttt{/proc} do Linux. O \texttt{/proc} é um sistema de arquivos virtual mantido pelo \textit{kernel} que expõe contadores acumulados de \textit{bytes} , pacotes, interrupções e estados de processos em arquivos de texto legíveis em espaço de usuário \cite{linux_proc_man}. Embora sua interface seja simples e portável, estudos demonstram que a leitura frequente de \texttt{/proc/net/dev} em ambientes com alto volume de tráfego pode gerar contenção na pilha de rede do \textit{kernel}, uma vez que os contadores são protegidos por \textit{spinlocks} e sua serialização para texto introduz latência de cópia proporcional ao número de interfaces monitoradas \cite{10.1145/3452296.3472888}.

O Sysstat se apresenta como uma ferramenta amplamente consolidada ao longo dos anos. Ela é um conjunto de utilitários que extrai estatísticas de desempenho lendo diretamente o \texttt{/proc} do Linux \cite{godard2024sysstat}. Utilitários como \texttt{sar} e \texttt{iostat} realizam leituras periódicas de arquivos como \texttt{/proc/net/dev}, \texttt{/proc/stat} e \texttt{/proc/[pid]/stat}, calculando deltas entre amostras consecutivas para estimar taxas de atividade. Essa abordagem se baseia em \textit{polling} periódico, gerando o ``efeito de observador'': sob alta frequência de coleta, o próprio ato de ler e processar os arquivos \texttt{/proc} consome ciclos significativos de CPU, o que pode atrasar o tempo de resposta da VNF monitorada \cite{10.1109/NOMS47738.2020.9110434}. Em ambientes com múltiplos processos sendo rastreados simultaneamente, a sobrecarga do \textit{polling} pode se acumular e se tornar não negligenciável, especialmente quando o intervalo de amostragem é inferior a um segundo \cite{10.1145/3452296.3472888}.

Outra ferramenta que vem ganhando espaço atualmente é o Prometheus. Consolidando-se como o padrão para observabilidade em arquiteturas de microsserviços e sistemas baseados em contêineres \cite{turnbull_prometheus_2018, usman_observability_2022}, seu modelo de coleta é baseado em \textit{pull}: um servidor central realiza requisições HTTP periódicas a \textit{exporters}, agentes que expõem métricas em um endpoint \texttt{/metrics} padronizado \cite{usman_observability_2022}. Esses exporters, como o \textit{Node Exporter}, coletam métricas do sistema operacional via \texttt{/proc} e as disponibilizam via HTTP para que o servidor Prometheus as consuma e armazene como séries temporais. Nesse modelo, a sobrecarga de monitoramento inclui tanto o \textit{polling} interno do exporter quanto o custo de serialização e transporte HTTP, o que pode introduzir atrasos na captura de eventos de curta duração em VNFs \cite{10.1109/NOMS47738.2020.9110434}. Na variante experimental deste trabalho, sem servidor central de coleta, o exporter opera de forma autônoma, mantendo o endpoint \texttt{/metrics} disponível para consultas eventuais. Nesse cenário, a sobrecarga adicional em relação ao Prometheus decorre da manutenção do servidor HTTP e da serialização contínua das métricas em memória.

\section{A Revolução da Programabilidade no \textit{kernel}: eBPF}
\label{sec:ebpf}

As origens do eBPF remontam ao trabalho de \citeonline{260309}, que propuseram em 1993 o \textit{BSD Packet Filter} (BPF) como uma arquitetura de máquina virtual embutida no \textit{kernel} para filtragem eficiente de pacotes de rede. O BPF original utilizava um conjunto de instruções minimalista com dois registradores e uma máquina de pilha, permitindo que filtros escritos em linguagem de alto nível fossem compilados para \textit{bytecode} e executados diretamente no \textit{kernel} sem cópia desnecessária de pacotes para o espaço do usuário. Essa abordagem eliminou a ineficiência das implementações anteriores, que exigiam que todos os pacotes fossem copiados para o espaço do usuário antes de qualquer filtragem, reduzindo significativamente a sobrecarga de captura de tráfego.

Ao longo das décadas seguintes, o BPF foi progressivamente estendido no \textit{kernel} Linux até que, em 2014, foi redesenhado como \textit{Extended BPF} (eBPF), com um conjunto de instruções de 64 bits, 11 registradores e suporte a mapas compartilhados entre \textit{kernel} e espaço do usuário. Essa extensão transformou o eBPF de um filtro de pacotes em um motor de programabilidade de propósito geral no \textit{kernel} \cite{gregg_bpf_2019, 10.1109/NOMS47738.2020.9110434}. O eBPF surge, assim, como uma tecnologia disruptiva capaz de estender as funcionalidades do \textit{kernel} do sistema operacional de forma dinâmica e segura, sem a necessidade de carregar módulos adicionais ou reiniciar o \textit{host}.

O eBPF opera por meio de uma máquina virtual embutida no \textit{kernel} Linux, na qual programas em \textit{bytecode} são validados por um verificador estático (\textit{verifier}) antes de serem carregados. Este verificador garante que o código não contenha laços infinitos e que todos os acessos à memória sejam seguros, evitando falhas catastróficas (\textit{kernel crashes}). Após a validação, os programas são compilados \textit{Just-In-Time} (JIT) para instruções nativas do processador, garantindo desempenho equivalente ao código compilado diretamente no \textit{kernel}. Para coletar métricas, os programas eBPF se anexam a ganchos (\textit{hooks}) do sistema \cite{gregg_bpf_2019, 10.1109/NOMS47738.2020.9110434, mizinski2025impact}, tais como:

\begin{itemize}
    \item \textbf{kprobes e uprobes:} Instrumentação dinâmica de funções do \textit{kernel} e do espaço do usuário, respectivamente. As \textit{kprobes} permitem associar um programa eBPF a qualquer ponto de entrada ou saída de função interna do \textit{kernel}, como \texttt{tcp\_sendmsg}, \texttt{tcp\_cleanup\_rbuf} ou \texttt{sock\_recvmsg}, interceptando parâmetros e valores de retorno sem modificar o código do \textit{kernel}. Essa característica as torna particularmente úteis para o monitoramento de comportamento de VNFs, pois permitem contabilizar bytes transmitidos e recebidos por conexão, medir a latência de chamadas de sistema e rastrear alocações de memória diretamente no ponto de execução, sem necessidade de copiar dados para o espaço do usuário. As \textit{uprobes}, por sua vez, operam de forma análoga, porém instrumentando funções em bibliotecas e executáveis no espaço do usuário, sendo úteis para rastrear o comportamento interno de aplicações sem acesso ao código-fonte;

    \item \textbf{Traffic Control (TC):} Inspeção e manipulação de pacotes diretamente nas filas do subsistema de rede do \textit{kernel}, nos pontos de entrada (\texttt{ingress}) e saída (\texttt{egress}) de cada interface de rede virtual ou física. Diferentemente do XDP, que opera antes da alocação da estrutura \texttt{sk\_buff}, os programas eBPF associados ao TC têm acesso ao pacote já encapsulado nessa estrutura, o que permite inspecionar e modificar campos de cabeçalhos de protocolos das camadas 3 e 4, como endereços IP, portas TCP/UDP e flags de controle, além de coletar métricas por fluxo \cite{10.1145/3594255.3594256}. O TC é amplamente utilizado em soluções de rede para contêineres e ambientes Kubernetes, como o Cilium, para implementar políticas de segurança, balanceamento de carga e observabilidade sem a necessidade de proxies no espaço do usuário. Em contextos de monitoramento de VNFs, o TC oferece um ponto de coleta privilegiado para medir a vazão e a latência de pacotes em interfaces virtuais, com sobrecarga inferior à de ferramentas baseadas em \texttt{/proc};

    \item \textbf{eXpress Data Path (XDP):} Processamento ultra-rápido de pacotes na camada mais baixa da pilha de rede do \textit{kernel}, executado diretamente no driver da placa de rede (\textit{Network Interface Card} - NIC), antes mesmo da alocação de estruturas complexas de memória como a \texttt{sk\_buff}. Ao operar nesse ponto, o XDP elimina grande parte da sobrecarga associada ao processamento convencional de pacotes no \textit{kernel}, permitindo taxas de processamento da ordem de dezenas de milhões de pacotes por segundo em hardware comum \cite{10.1145/3594255.3594256}. Programas XDP podem tomar decisões de encaminhamento (\texttt{XDP\_PASS}, \texttt{XDP\_DROP}, \texttt{XDP\_TX}, \texttt{XDP\_REDIRECT}) e coletar métricas de tráfego com latência mínima, tornando-o adequado para casos de uso como mitigação de ataques DDoS, balanceamento de carga de alto desempenho e coleta de telemetria em redes de alta velocidade \cite{technologies12080122, 10.1145/3594255.3594256}. A principal limitação do XDP em relação ao TC é o acesso restrito ao contexto do pacote: por ser executado antes da alocação da \texttt{sk\_buff}, o programa XDP não tem acesso a informações de estado de conexão mantidas pelo \textit{kernel}, como o estado da máquina de estados TCP, o que limita sua aplicabilidade para monitoramento de métricas orientadas a fluxo \cite{gregg_bpf_2019, shahinfar2025demystifying}.
\end{itemize}

Para a troca de informações entre o \textit{kernel} e o espaço do usuário, o eBPF utiliza os mapas eBPF (BPF \textit{Maps}), que são estruturas chave-valor compartilhadas de forma assíncrona. Essa arquitetura permite realizar agregações estatísticas complexas diretamente no \textit{kernel}, de modo que o monitor no espaço do usuário precise ler apenas as métricas consolidadas periodicamente, eliminando custos excessivos de comunicação e alternância de contexto \cite{gregg_bpf_2019}. A portabilidade dos programas eBPF entre versões de \textit{kernel} é viabilizada pela abordagem CO-RE (\textit{Compile Once, Run Everywhere}), implementada pela biblioteca \textit{libbpf}. O CO-RE utiliza metadados de tipos de dados embutidos no \textit{kernel} no formato BTF (BPF \textit{Type Format}) para relocar automaticamente, em tempo de carregamento, os acessos a campos de estruturas internas do \textit{kernel} que possam ter mudado de posição entre versões \cite{nakryiko2020core}. Essa técnica elimina a dependência de ferramentas de compilação em tempo de execução (como o BCC, que exigia cabeçalhos de \textit{kernel} e LLVM no \textit{host}), permitindo que o \textit{bytecode} compilado uma única vez seja carregado em \textit{kernels} distintos sem recompilação \cite{nakryiko2020core, gregg_bpf_2019}.

Apesar de ser amplamente considerado o estado da arte em observabilidade, o eBPF não é uma solução isenta de limitações. Pesquisas recentes, \cite{shahinfar2025demystifying}, demonstraram que o uso incorreto de eBPF pode, paradoxalmente, degradar o desempenho de aplicações de rede.

A Arquitetura de Conjunto de Instruções (\textit{Instruction Set Architecture} - ISA) do eBPF não oferece suporte nativo para operações complexas, como aritmética de ponto flutuante ou instruções vetoriais SIMD. Além disso, o compilador JIT do \textit{kernel} frequentemente gera instruções de máquina menos otimizadas do que o compilador do espaço do usuário (por exemplo, em operações de cópia de memória, o eBPF pode se comportar de forma até dez vezes mais lenta do que um módulo de \textit{kernel} nativo devido à falta de instruções nativas otimizadas). Outro ponto crítico refere-se à violação do isolamento de desempenho (\textit{performance isolation}): como os programas eBPF são executados antes que ocorra a multiplexação das conexões do sistema, um programa de monitoramento ineficiente associado a uma determinada VNF pode interferir diretamente na latência e na vazão de outras aplicações hospedadas no mesmo servidor \cite{shahinfar2025demystifying}.

Sob carga, no entanto, soluções de rede e monitoramento baseadas em eBPF superam os modelos tradicionais baseados em \textit{iptables} \cite{mizinski2025impact}. Em clusters de contêineres sob estresse, o uso do eBPF reduz significativamente a oscilação de latência e diminui o consumo geral de CPU do sistema de 45,83\% para 40,94\%, evidenciando sua robustez para ambientes virtualizados de alta densidade \cite{mizinski2025impact}.

\section{Monitoramento do Comportamento de VNFs: \textit{kprobes} versus \texttt{/proc}}
\label{sec:kprobes_proc}

No contexto de monitoramento de VNFs, duas abordagens fundamentais se contrapõem para a coleta de métricas de tráfego: a instrumentação dinâmica no \textit{kernel} por meio de \textit{kprobes} e a leitura periódica do pseudossistema de arquivos \texttt{/proc} no espaço do usuário.
As \textit{kprobes} permitem que programas eBPF sejam associados a pontos de entrada e saída de funções internas do \textit{kernel}, como \texttt{tcp\_sendmsg} e \texttt{tcp\_cleanup\_rbuf}, que são invocadas a cada operação de envio e recepção de dados TCP, respectivamente \cite{gregg_bpf_2019}. Ao interceptar essas funções, é possível acumular contadores de bytes transmitidos e recebidos por porta diretamente no \textit{kernel}, sem cópias de pacotes para o espaço do usuário e sem pressão adicional sobre a pilha de rede \cite{10.1109/NOMS47738.2020.9110434}.

Os resultados são disponibilizados ao espaço do usuário de forma assíncrona por meio de mapas eBPF, reduzindo a alternância de contexto e os custos de comunicação inter-espaço \cite{gregg_bpf_2019}.
Para ilustrar o impacto dessa abordagem em um ambiente NFV, considere um WAF virtualizado submetido a um tráfego de 50.000 requisições HTTP por segundo. Um monitor baseado em \textit{kprobes} associado a \texttt{tcp\_sendmsg} acumula os contadores de bytes por conexão diretamente no \textit{kernel} e os disponibiliza ao espaço do usuário apenas quando consultados, sem interferir no caminho crítico de processamento das requisições \cite{10.1109/NOMS47738.2020.9110434}. O MANO pode, assim, consultar o mapa eBPF a cada segundo e obter métricas de vazão por porta com granularidade de conexão, sem que o processo de coleta consuma ciclos de CPU do próprio WAF.

Em um cenário de SFC composta por \textit{firewall}, IDS e balanceador de carga, essa abordagem permite monitorar simultaneamente cada elo da cadeia com sobrecarga aditiva mínima, pois cada programa eBPF opera de forma independente e assíncrona no \textit{kernel} \cite{kannan2024designing}.
Em contrapartida, o pseudossistema de arquivos \texttt{/proc/net/dev} disponibiliza contadores acumulados de bytes e pacotes por interface de rede, mantidos pelo próprio \textit{kernel}. Ferramentas como o Sysstat e o Prometheus realizam leituras periódicas desses contadores, calculando deltas entre amostras consecutivas para estimar a taxa de tráfego \cite{godard2024sysstat}.

Embora esse mecanismo seja portável e de baixa complexidade de implementação, a frequência de amostragem e o processamento adicional no espaço do usuário introduzem ciclos de CPU concorrentes ao processo monitorado, podendo elevar o tempo de resposta da VNF \cite{10.1109/NOMS47738.2020.9110434}. Em particular, estudos sobre o custo de instrumentação da pilha de rede do Linux mostram que leituras síncronas de contadores de rede protegidos por \textit{spinlocks} podem causar contenção sob alto volume de tráfego, aumentando o tempo médio de resposta de conexões concorrentes \cite{10.1145/3452296.3472888}.

Para ilustrar o impacto em um ambiente NFV, considere o mesmo WAF virtualizado operando sob carga equivalente. Um agente Sysstat configurado para amostrar \texttt{/proc/net/dev} a cada segundo executa, em cada ciclo, a leitura e a serialização dos contadores de todas as interfaces de rede do host, incluindo interfaces virtuais de outros contêineres ou VMs co-localizados. Sob alta densidade de VNFs no mesmo servidor físico, esse custo se multiplica proporcionalmente ao número de interfaces monitoradas, podendo introduzir atrasos perceptíveis no tempo de resposta do WAF \cite{10.1145/3452296.3472888}. Em um cenário de SFC, o impacto é ainda mais crítico: se cada VNF da cadeia executa seu próprio agente de \textit{polling}, a sobrecarga total no host pode se tornar expressiva, comprometendo a capacidade de reação do MANO a eventos de curta duração, como picos abruptos de tráfego ou falhas transitórias em um dos elos \cite{oliveira_observability_2026, 10.23919}. Adicionalmente, a granularidade temporal limitada pelo intervalo de amostragem impede a detecção de anomalias de curta duração: um pico de latência que dura menos de um segundo é completamente mascarado pela média calculada entre duas amostras consecutivas \cite{epping2024rtt}.

Diante do exposto, as três abordagens apresentam \textit{trade-offs} distintos com impactos técnicos diretos sobre o desempenho das VNFs monitoradas. O Sysstat, ao realizar \textit{polling} periódico via \texttt{/proc}, executa no mesmo espaço de endereçamento do sistema operacional que hospeda a VNF, competindo diretamente por ciclos de CPU com o processo monitorado. Sob intervalos de amostragem inferiores a um segundo, esse comportamento pode se manifestar como aumento mensurável no tempo de resposta da VNF, especialmente em cenários de alta taxa de requisições, nos quais cada ciclo de CPU desviado para leitura e serialização de \texttt{/proc} representa capacidade de processamento subtraída do tráfego de serviço.

O Prometheus, por sua vez, acrescenta à sobrecarga do \textit{polling} interno do \textit{exporter} o custo de manutenção de um servidor HTTP e da serialização contínua das métricas em memória; em VNFs de baixa folga computacional, como um WAF operando próximo ao limite de capacidade, essa sobrecarga adicional pode ser suficiente para elevar a cauda de latência das requisições além do limiar de SLA. Ambas as abordagens compartilham uma limitação estrutural: por operarem com granularidade temporal limitada pelo intervalo de amostragem, são incapazes de capturar eventos transitórios de curta duração, como picos de latência em conexões individuais ou rajadas de tráfego subsegundo, que podem ser sintomas precoces de degradação funcional da VNF.

O eBPF, ao instrumentar o \textit{kernel} diretamente por meio de \textit{kprobes}, elimina o \textit{polling} no espaço do usuário e reduz a alternância de contexto entre kernel e aplicação, uma vez que as agregações estatísticas são realizadas no próprio \textit{kernel} e apenas os resultados consolidados são transferidos ao espaço do usuário via mapas eBPF. Esse modelo reduz tanto o número de transições de contexto quanto o volume de dados transferidos entre os dois espaços, o que se traduz em menor interferência sobre o caminho crítico de processamento da VNF. Contudo, o eBPF não é isento de impacto: cada ponto de instrumentação introduz uma chamada de função adicional no caminho de execução do \textit{kernel}, e programas mal dimensionados podem degradar o desempenho de todas as aplicações hospedadas no mesmo servidor devido à violação do isolamento de desempenho. Em termos práticos, para uma VNF como o WAF, o impacto do eBPF sobre o tempo de resposta tende a ser significativamente menor que o das abordagens baseadas em \texttt{/proc} sob alta carga, conforme evidenciado pela redução de consumo de CPU de 45,83\% para 40,94\% reportada por \citeonline{mizinski2025impact} em ambientes Kubernetes sob estresse. Embora a literatura indique essas vantagens do eBPF, há escassez de estudos quantitativos controlados que mensurem diretamente a sobrecarga dessas três abordagens sobre uma mesma VNF em condições equivalentes \cite{usman_observability_2022}, lacuna que o presente trabalho busca endereçar empiricamente.


\section{Trabalhos Relacionados}
\label{sec:trabalhos_relacionados}

Diversos trabalhos investigaram o uso de eBPF para monitoramento de rede e a questão da observabilidade em ambientes virtualizados, porém com escopos, abordagens e métricas distintas das adotadas neste trabalho.

\citeonline{10.1109/NOMS47738.2020.9110434} realizam uma análise qualitativa e experimental do eBPF como alternativa não intrusiva para monitoramento de desempenho, contrastando-o com ferramentas tradicionais baseadas em \textit{polling} no espaço do usuário. Os autores demonstram que programas eBPF associados a \textit{kprobes} introduzem latência de instrumentação significativamente menor do que abordagens baseadas em \texttt{/proc}, especialmente sob alta taxa de eventos de rede. Contudo, o trabalho não considera o Prometheus como alternativa de comparação nem quantifica a sobrecarga sobre uma VNF específica sob cargas controladas e crescentes.

\citeonline{kannan2024designing} propõem um agente de observabilidade de rede leve (\textit{netobserv\discretionary{-}{}{}ebpf\discretionary{-}{}{}agent}) construído inteiramente sobre eBPF, voltado para ambientes Kubernetes. O agente coleta métricas de fluxo TCP/UDP diretamente no \textit{kernel}, sem passar pelo espaço do usuário, obtendo redução no consumo de CPU e memória em relação a agentes tradicionais. A principal diferença em relação ao presente trabalho é que o foco recai sobre a coleta de metadados de fluxo em larga escala, não sobre a quantificação do impacto do monitoramento no tempo de resposta de uma VNF de segurança.

\citeonline{epping2024rtt} investigam o monitoramento passivo de latência TCP no \textit{kernel} Linux utilizando eBPF, implementando o estimador \textit{epping} para medir RTT de conexões ativas sem alterar o tráfego. O trabalho evidencia a viabilidade de coletar métricas de desempenho de rede com overhead mínimo diretamente no \textit{kernel}, corroborando a hipótese deste trabalho sobre as vantagens de abordagens \textit{kernel-level}. Diferentemente deste estudo, os autores não comparam o eBPF com ferramentas de espaço de usuário sob condições de carga equivalentes.

\citeonline{shahinfar2025demystifying} desmistificam o desempenho de aplicações de rede baseadas em eBPF, demonstrando que o uso inadequado de programas eBPF pode degradar o desempenho de aplicações hospedadas no mesmo servidor devido à violação do isolamento de desempenho. Embora não comparem eBPF com Sysstat ou Prometheus, os resultados fundamentam a importância de quantificar experimentalmente a sobrecarga de cada abordagem, conforme proposto neste trabalho.

\citeonline{10.1145/3594255.3594256} investigam o uso de estruturas de dados probabilísticas (\textit{sketches}) implementadas em eBPF para monitoramento de tráfego em alta velocidade, demonstrando que é possível realizar medições de fluxo com precisão controlável diretamente no plano de dados do \textit{kernel}. O trabalho aponta que a abordagem elimina a necessidade de exportar pacotes individuais para o espaço do usuário, reduzindo drasticamente o volume de dados transmitidos. Embora o contexto seja distinto, focado em telemetria de fluxo em redes de alta velocidade em vez de VNFs individuais, os resultados corroboram a viabilidade de agregações eficientes no \textit{kernel}, que é a mesma técnica empregada pelo observador eBPF deste trabalho.

Do ponto de vista da observabilidade em ambientes NFV, \citeonline{10.23919} discutem os desafios de observabilidade para VNFs em sistemas 5G, argumentando que ferramentas existentes carecem de granularidade e baixa intrusividade para ambientes de alta densidade de funções virtualizadas. Os autores identificam como lacuna a ausência de mecanismos de coleta que não comprometam o desempenho da própria VNF monitorada, lacuna que este trabalho aborda empiricamente.

Em perspectiva mais recente, \citeonline{oliveira_observability_2026} realizam uma análise conceitual do papel da observabilidade no gerenciamento de ambientes NFV, propondo um modelo teórico que relaciona métricas de monitoramento, latência de coleta e capacidade de reação do orquestrador. Os autores conduzem experimentos quantitativos em um ambiente controlado com uma VNF que simula um DPI, demonstrando que aumentos na latência de coleta de métricas comprometem diretamente a eficácia de decisões automatizadas de escalonamento e recuperação. Contudo, o estudo se limita a uma única abordagem de monitoramento com o Sysstat (\textit{polling} via \texttt{/proc}) e não contempla a comparação quantitativa entre ferramentas concorrentes como Prometheus e eBPF.


O presente trabalho tem como objetivo suprir essa lacuna ao comparar experimentalmente a sobrecarga imposta por essas três tecnologias de monitoramento, eBPF, Sysstat e Prometheus, sobre o tempo de resposta de uma VNF sob diferentes volumes de carga. Para sintetizar as características de cada iniciativa e evidenciar a lacuna preenchida por esta pesquisa, a Tabela~\ref{tab:trabalhos_relacionados_final} consolida a comparação entre os trabalhos relacionados e a abordagem proposta.

\begin{table}[htbp]
    \centering
    \caption{Análise comparativa entre os trabalhos relacionados e a presente pesquisa.}
    \label{tab:trabalhos_relacionados_final}
    \renewcommand{\arraystretch}{1.3} % Aumenta levemente o espaçamento entre as linhas
    \resizebox{\textwidth}{!}{ % Ajusta a tabela à largura da página
    \begin{tabular}{p{3cm} p{4cm} p{3.5cm} p{3cm} p{2.5cm} p{2.5cm}}
        \toprule
        \textbf{Referência} & \textbf{Foco Principal} & \textbf{Tecnologias Analisadas} & \textbf{Métricas de Sobrecarga} & \textbf{Carga Controlada} & \textbf{Contexto} \\
        \midrule

        Mijumbi et al. (2016) & Revisão dos desafios e estado da arte em NFV & N/A (Teórico) & N/A & Não & NFV Geral \\

        Cassagnes et al. (2020) & Comparativo entre eBPF e polling em espaço de usuário & eBPF (kprobes) e ferramentas tradicionais & Latência de instrumentação & Sim (Alta taxa de rede) & Rede Geral / Host \\

        Cai et al. (2021) & Overhead da pilha de rede do Linux & Ferramentas baseadas em \textit{/proc} & CPU, Contenção (\textit{Spinlocks}) & Sim (Alto tráfego) & Host (Bare-metal) \\

        Usman et al. (2022) & \textit{Survey} sobre observabilidade em microsserviços & Prometheus, Sysstat (citados no modelo \textit{pull}) & N/A & Não & Microsserviços / \textit{Edge} \\

        Miano et al. (2023) & Telemetria e medição de fluxo em alta velocidade & eBPF (TC / XDP) & Precisão vs. Overhead & Sim & Plano de dados do Kernel \\

        Sathyaseelan et al. (2024) & Criação de agente de rede leve para observabilidade & eBPF & CPU, Memória & Sim & Cluster Kubernetes \\

        Shahinfar et al. (2025) & Degradação de performance e isolamento no eBPF & eBPF & \textit{Throughput}, Latência & Sim (Concorrência) & Host / Rede \\

        Oliveira et al. (2026) & Impacto da latência de coleta na orquestração & Sysstat (\textit{/proc}) & Latência de coleta, Reação do MANO & Sim (DPI simulado) & NFV / Orquestração \\

        \textbf{Este trabalho} & \textbf{Benchmarking de impacto de ferramentas na VNF} & \textbf{eBPF, Sysstat, Prometheus} & \textbf{CPU, RAM, RTT UDP (ns)} & \textbf{Sim (100k a 2M msgs)} & \textbf{VNF (WAF Isolado)} \\
        \bottomrule
    \end{tabular}
    }
\end{table}
